\begin{definition}
	Функциональный ряд $\ss f_n(x)$ \textit{сходится равномерно} на множестве $E$, если на $E$ равномерно сходится последовательность его частичных сумм $\{\sum_{k=1}^nf_k(x)\}_{n=1}^\i$.
\end{definition}
\begin{theorem}\nf{(Критерий Коши равномерной сходимости функциональных рядов)}\label{15_th1}\\
	Функциональный ряд $\ss f_n(x)$ сходится равномерно на $E$ тогда и только тогда, когда:
	\[(\fa\eps>0)(\ex N\in\N)(\fa n>N)(\fa p\in\N)(\fa x\in E)\ \l|\sum_{k=n+1}^{n+p}f_k(x)\r|<\eps\]
\end{theorem}
\begin{anote}
	\textbf{Доказательство.} $\Ra$: Пусть $\lim_{n\ra\i}\sum_{k=1}^nf_n(x)=f(x)$.
	\begin{equation*}
		\begin{split}
			(\fa\eps>0)(\ex N\in\N)(\fa n>N)(\fa x\in E&)\ \l|\sum_{k=1}^nf_k(x)-f(x)\r|<\frac\eps2\\
			(\fa p\in\N&)\ \l|\sum_{k=1}^{n+p}f_k(x)-f(x)\r|<\frac\eps2
		\end{split}
	\end{equation*}
	Получили:
	\[(\fa\eps>0)(\ex N\in\N)(\fa n>N)(\fa x\in E)(\fa p\in\N)\ \l|\sum_{k=n+1}^{n+p}f_k(x)\r|<\eps\]
	$\La$: По признаку Коши сходимости функциональных последовательностей имеем:
	\[\lim_{n\ra\i}\sum_{k=1}^nf_k(x)=f(x)\]
	Тогда в имеющемся неравенстве устремим $p\ra\i$, тогда:
	\[(\fa\eps>0)(\ex N\in\N)(\fa n>N)(\fa x\in E)\ \l|\sum_{k=1}^\i f_k(x)-f(x)\r|<\eps\]
	Получили требуемое.
\end{anote}
\begin{corollary}\nf{(Необходимое условие равномерной сходимости функциональных рядов)}\\
	Если функциональный ряд сходится равномерно на $E$, то $f_n\rra0$ на $E$.
\end{corollary}
\begin{proof}
	Возьмём в доказательстве (\ref{15_th1}) $p=1$, тогда получим требуемое.
\end{proof}
\begin{note}
	В следствии (\ref{14_cr1}) говорится:
	\[f_n\text{ не сходится равномерно к }f\text{ на }E\lra(\ex\{x_n\}\subset E)\ f_n(x_n)-f(x_n)\nra0 \text{ при }n\ra\i\]
	Из того, что $\ss f_n(x_n)$ расходится для некоторой последовательности не следует, что $\ss f_n(x)$ не сходится равномерно на $E$.
\end{note}
\begin{theorem}\nf{(Признак Вейерштрасса, M-test)}\label{15_th2}\\
	Если $(\fa x\in E)(\fa n\in\N)\ |f_n(x)|\le c_n$ и $\ss c_n$ сходится, то $\ss f_n(x)$ сходится равномерно на $E$.
\end{theorem}
\begin{proof}
	\[\ss c_n\text{ сходится }\Ra(\fa\eps>0)(\ex N\in\N)(\fa n>N)(\fa p\in\N)\ \sum_{k=n+1}^{n+p}c_k<\eps\]
	Тогда имеем:
	\[\l|\sum_{k=n+1}^{n+p}f_k(x)\r|\le\sum_{k=n+1}^{n+p}|f_k(x)|\le\sum_{k=n+1}^{n+p}c_k<\eps\]
\end{proof}
\begin{definition}
	Назовём функциональную последовательность $\sq[f_n(x)]$ \textit{равномерно ограниченной} на $E$, если:
	\[(\ex c)(\fa x\in E)(\fa n\in\N)\ |f_n(x)|\le c\]
\end{definition}
\begin{theorem}\nf{(Признак Дирихле равномерной сходимости функциональных рядов)}\label{15_th3}\\
	Если:
	\begin{enumerate}
		\item Частичная сумма $\sum_{k=1}^nf_k(x)$ равномерно ограничена на $E$.
		\item $g_n(x)$ монотонна $(\fa x\in E)$ и $g_n\rra0$ на $E$.
	\end{enumerate}
	то $\ss f_n(x)g_n(x)$ равномерно сходится на $E$.
\end{theorem}
\begin{proof}
	Хотим доказать:
	\[(\fa\eps>0)(\ex N\in\N)(\fa n>N)(\fa p\in\N)(\fa x\in E)\ \l|\sum_{k=n+1}^{n+p}f_k(x)g_k(x)\r|<\eps\]
	Пусть $F_i(x):=\sum_{k=n+1}^if_K(x),\ F_n(x)=0$.\\
	\begin{equation*}
		\begin{split}
			\l|\sum_{k=n+1}^{n+p}f_k(x)g_k(x)\r|&=
			\sum_{k=n+1}^{n+p}\Big(F_k-F_{k-1}(x)\Big)\cdot\Big|g_k(x)\Big|=\\
			&=\l|F_{n+p}(x)\cdot g_{n+p}+\sum_{k=n+1}^{n+p-1}F_k(x)\cdot\Big(g_k(x)-g_{k+1}(x)\Big)\r|\le\\
			&\le2\cdot c\cdot\l(\Big|g_{n+p}(x)\Big|+\sum_{k=n+1}^{n+p}\Big|g_k(x)-g_{k+1}(x)\Big|\r)=\\
			&=2\cdot c\cdot\l(\bigg|g_{n+p}(x)\bigg|+\l|\sum_{k=n+1}^{n+p}\Big(g_k(x)-g_{k+1}(x)\Big)\r|\r)\le\\
			&\le2\cdot c\cdot\Big(2\cdot|g_{n+p}(x)|+|g_{n+1}(x)|\Big)<\eps
		\end{split}
	\end{equation*}
	Получили требуемое.
\end{proof}
\begin{theorem}\nf{(Признак Абеля равномерной сходимости функциональных рядов)}\label{15_th4}\\
	Если:
	\begin{enumerate}
		\item $\ss f_n(x)$ равномерной сходится на $E$.
		\item $\sq[g_n(x)]$ монотонна и равномерно ограничена на $E$.
	\end{enumerate}
	То $\ss f_n(x)g_n(x)$ равномерно сходится на $E$.
\end{theorem}
\begin{proof}
	Пусть $M:=(\fa x\in E)(\fa n\in\N)\ |g_n(x)|<M$
	По (\ref{15_th1}):
	\[(\fa\eps>0)(\ex N\in\N)(\fa n>N)(\fa p\in\N)(\fa x\in E)\ \l|\sum_{k=n+1}^{n+p}f_k(x)\r|<\frac\eps{6M}\]
	Перепишем часть доказательства из теоремы выше:
	\begin{equation*}
		\begin{split}
			\l|\sum_{k=n+1}^{n+p}f_k(x)g_k(x)\r|&=
			\sum_{k=n+1}^{n+p}\Big(F_k-F_{k-1}(x)\Big)\cdot\Big|g_k(x)\Big|=\\
			&=\l|F_{n+p}(x)\cdot g_{n+p}+\sum_{k=n+1}^{n+p-1}F_k(x)\cdot\Big(g_k(x)-g_{k+1}(x)\Big)\r|\le\\
			&\le\frac\eps{6M}\l(\Big|g_{n+p}(x)\Big|+\sum_{k=n+1}^{n+p}\Big|g_k(x)-g_{k+1}(x)\Big|\r)\le\eps
		\end{split}
	\end{equation*}
	Получили требуемое.
\end{proof}
\begin{example}
	Рассмотрим ряд $\ds\ss\frac{\sin nx}{n^\al}$. При $\al\le0$ ряд расходится.\\
	\[\al>1:\ \l|\frac{\sin x}{n^\al}\r|\le\frac1{n^\al},\ \ss\frac1{n^\al}\text{ сходится }\Ra\ss\frac{\sin nx}{n^\al}\text{ сходится равномерно на }\R\]
	\[0<\al\le1:\ \l|\sum_{k=1}^n\sin x\r|\le\frac2{|\sin\frac x2|},\ x\ne2k\pi,\ k\in\Z\]
	\[\frac2{|\sin\frac x2|}\le c\text{ на }[a,b]\subset(0,2\pi)\]
	При $0<\al\le1$ ряд сходится равномерно на $\fa[a,b]$. На $[0,\delta]$ ряд не сходится равномерно.\\
	Запишем отрицание критерия Коши:
	\[(\ex\eps>0)(\fa N\in\N)(\ex n>N)(\ex p\in\N)(\ex x\in[0,\delta])\ \l|\sum_{k=n+1}^{n+p}\frac{\sin nx}{n^\al}\r|\ge\eps\]
	Пусть $p=n$, тогда $\frac\pi3<kx\le\frac{2\pi}3$. В итоге имеем:
	\[\sum_{k=n+1}^{2n}\frac{\sin kx}{k^\al}\ge\frac{\sqrt{3}}2\cdot\frac1{2n}\cdot n=\frac{\sqrt{3}}4\]
	Получили отрицание критерия Коши.
\end{example}
\begin{example}
	Рассмотрим $\sq[x^n]$ на $[0,1]$.
	\[\lim_{n\ra\i}x^n=
	\begin{cases}
		0,\ x\in[0,1)\\
		1,\ x=1
	\end{cases}
	\]
	Какое бы большое $n$ мы ни выбирали, с какого - то момента последовательность будет стремиться к 1.
\end{example}
\begin{theorem}\nf{(Предельный переход в функциональных последовательностях)}\label{15_th5}\\
	Пусть $\sq[f_n(x)]$ равномерно сходится к $f(x)$ на метрическом пространстве $E$, $x_0$ - предельная точка, $\ex\lim_{x\ra x_0,x\in E}f_n(x)=a_n\in\R$. Тогда $\ex\lim_{n\ra\i}a_n=\lim_{x\ra x_0}f(x)$.
\end{theorem}
\begin{proof}
	По критерию Коши:
	\[(\fa\eps>0)(\ex N\in\N)(\fa n>N)(\fa p\in\N)(\fa x\in E)\ |f_{n+p}(x)-f_n(x)|<\eps\]
	Перепишем при $x\ra x_0$:
	\[|a_{n+p}-a_n|\le\eps\Ra\ex\lim_{n\ra\i}a_n=a\]
	\[|f(x)-a|\le|f(x)-f_n(x)|+|f_n(x)-a_n|+|a_n-a|\]
	Теперь оценим полученное выражение.\\
	По условию имеем:
	\[(\ex N\in\N)(\fa n>N)(\fa x\in E)\ |f(x)-f_n(x)|<\frac\eps3,\ |a_n-a|<\frac\eps3\]
	\[(\ex\delta>0)(\fa x\in E:0<\rho(x,x_n))<\delta)\ |f_n(x)-a_n|<\frac\eps3\]
	Итого имеем:
	\[(\fa\eps>0)(\ex N\in\N)(\fa n>N)(\fa p\in\N)(\fa x\in E)(\ex\delta>0:0<\rho(x,x_n)<\delta)\ |f(x)-a|<\eps\]
\end{proof}
\begin{corollary}
	Если $(\fa n\in\N)\ f_n$ непрерывна и равномерно сходится к $f$ на $E$, то $f$ непрерывна на $E$.
\end{corollary}
\begin{theorem}\nf{(Признак Дини равномерной сходимости функциональной последовательности)}\label{15_th6}\\
	Пусть $E$ - компактное метрическое пространство. Если $(\fa x\in E)\sq[f_n(x)]$ - убывающая последовательность непрерывных на $E$ функций, сходящаяся к непрерывной на $E$ функции, то эта сходимость равномерна.
\end{theorem}
\begin{proof}
	Положим: $f(x):=\lim_{n\ra\i}f_n(x),\ g_n(x):=f_n(x)-f(x)$.\\
	При $n\ra\i$ $g_n(x)$ убывает и стремится к 0. Так как $g_n$ непрерывна, то:
	\[(\fa\eps>0)(\fa x\in E)(\ex n_x\in\N)(\fa y\in U_\delta(x)\cap E)(\fa n\ge n_x)\ 0\le g_{n}(y)<\eps\]
	Получили открытое подпокрытие компактного множества.
	\[\bigcup_{x\in E}U_{\delta_x}(x)\supset E\Ra\ex\{x_1,\dots,x_N\}\subset E\text{ такое, что: }\bigcup_{i=1}^NU_{\delta_{x_i}}(x_i)\supset E\]
	\[(\fa x\in E)(\ex i\in\{1,\dots,N\})\ x\in U_{\delta_{x_i}}(x_i)\]
	Тогда имеем:
	\[(\fa x\in E)(\fa\eps>0)(\ex M=\max(n_{x_1},\dots, n_{x_N}))(\fa n>M)\ |f_n(x)-f(x)|<\eps\]
	Получили требуемое.
\end{proof}